\section{Lecture 2: 23rd January}

I discussed Thursday's lecture with Elliot Bobrow.

It seemed like the theme of the lecture was \emph{groups important to representation theory} --- the general linear group (primarily over $\mathbb{R}$), transformation groups, and symmetric groups.

For both of us, the main new idea we learned was that of a transformation group. I thought it was interesting that category theory could really help clarify what we mean by ``a transformation group is a subgroup of the automorphism group of a set''. Typically the case we're interested in is that of a transformation group $G = \operatorname{Aut}_{\mathscr{C}} X \le \operatorname{Aut}_{\mathsf{Set}} X$ for some (concrete) category $\mathscr{C}$ with a forgetful functor to $\mathsf{Set}$.

We were also both amazed that Cayley's theorem extends to infinite groups --- at the end Prof. Andr\'es mentioned that any group is a transformation group. Later we realised that this is just because $G \le \operatorname{Aut}_{\mathsf{Set}} G$ with the embedding that arises from $G$ acting on itself from the left.
